\chapter{Python programs for numerical computation}
\label{appendix-python-numerical}
\chapteroverview{The programs in this appendix connect transform identities
with working numerical algorithms. They implement fast transforms, integer
products, and Fourier solvers for differential equations. Explicit conventions
and independent reference calculations make it possible to understand both
the algorithms and the limits of floating-point arithmetic.}

\section*{Using the programs}

The printed code is intended to be run alongside the mathematical
derivations. This section explains how to use the accompanying modules and
check a first transform against an independent implementation.

The accompanying \texttt{listings/} directory contains ordinary Python
modules with their imports and shared helpers. The printed functions are
extracted directly from these files; functions in the same module can call
one another without being copied into separate files. From the source
directory, this example makes the modules available and checks an FFT:
\begin{lstlisting}[language=Python]
import sys
sys.path.insert(0, "listings")
import numpy as np
from fft_algorithms import fft_dit
x = np.array([1, 2, 0, -1], dtype=complex)
assert np.allclose(fft_dit(x), np.fft.fft(x))
\end{lstlisting}
The repository README gives installation and verification instructions.
NumPy arrays use zero-based indices, \texttt{a*b} is entrywise multiplication,
and \texttt{a @ b} is matrix multiplication. Numerical Fourier transforms
use the negative exponential and an unnormalized forward transform; the
inverse includes $1/N$. Unless otherwise stated, vectors are one-dimensional
arrays, polynomial coefficients increase with degree, and real data may
produce complex intermediate arrays. All random demonstrations fix their seed.

The implementations favor readable algorithms. Recursive versions allocate
temporary arrays; they are not in-place routines. For large numerical
computations, NumPy and SciPy provide compiled transforms; FFTW is another
established implementation \cite{fftw-software}. The checks supplied here
compare the teaching algorithms with direct mathematical formulas and
independent numerical or exact references.

\section{Fast Walsh transform}
\label{sect1-listing-fwt}

The Walsh recursion provides a compact example of a transform built entirely
from additions and subtractions. The implementation makes the tensor
factorization visible and keeps the inverse normalization explicit.

The function \texttt{fwt}, in \texttt{transforms.py}, follows the tensor
factorization of Section~\ref{sect2-fwt}. Its length must be a positive
power of two. The transform satisfies $H_N^2=NI$, so applying the same
function and dividing by $N$ gives its inverse. The displayed recursion
uses $O(N\log N)$ arithmetic and creates temporary arrays at each level.
\begin{listing}[H]
\PythonCode{listing-fwt}
\caption{A recursive fast Walsh transform.}
\label{listing-fwt}
\end{listing}

\section{Fast Hartley transform}
\label{sect1-listing-fht}

The fast Hartley transform shows how paired frequencies allow a real-kernel
transform to be evaluated recursively. The accompanying convolution routine
demonstrates how those coefficients can be used directly in a calculation.

The module \texttt{hartley.py} implements the real-kernel transform
of Section~\ref{sect2-algo-de-calcul-hartley}. The radix-two function
\texttt{fht} calls \texttt{hartley\_twiddle}, which combines the paired
frequencies $k$ and $-k$. Its inverse is again obtained by dividing
the forward transform by $N$. For real input, this recursion uses real
arithmetic throughout.

The convolution function accepts two nonempty vectors of the same length.
It uses \texttt{hartley}, an FFT-based helper valid at arbitrary lengths,
and combines the transforms of \emph{both} inputs. The result is circular
convolution. To compute linear convolution, first pad both vectors to a
common length at least the sum of their lengths minus one.
\begin{listing}[H]
\PythonCode{listing-fht}
\caption{A radix-two fast Hartley transform.}
\label{listing-fht}
\end{listing}

\begin{listing}[H]
\PythonCode{listing-operator_chi}
\caption{The paired-frequency twiddle operator for the FHT.}
\label{listing-operator_chi}
\end{listing}

\begin{listing}[H]
\PythonCode{listing-fht_convol}
\caption{Circular convolution from Hartley transforms of both inputs.}
\label{listing-fht_convol}
\end{listing}

\section{Fast Fourier transforms}
\label{sect1-listing-fft}

Different FFT factorizations compute the same transform but organize the
data differently. These implementations place recursive, time-decimation,
and frequency-decimation algorithms side by side so their indexing can be
compared.

The module \texttt{fft\_algorithms.py} provides three radix-two FFTs:
\texttt{fft\_recursive} follows the time-decimation identity directly;
\texttt{fft\_dit} iterates butterflies after bit reversal;
\texttt{fft\_dif} performs frequency-decimation butterflies before the
final reordering. All return frequencies in the usual order and accept
\texttt{inverse=True}, which reverses the phase and includes $1/N$.
Their mathematical derivations are in Sections~\ref{sect2-present-algo-fft},
\ref{sect2-impl-concrete}, and~\ref{sect2-frequency-decimation}.

The remaining helpers multiply by a linear phase, construct the full
bit-reversal permutation, and reverse the specified number of bits of
a single integer. The two iterative FFTs keep a working array and use
temporary vectors for each butterfly block; no claim of strict in-place
operation is needed.
\begin{listing}[H]
\PythonCode{listing-fft-rec}
\caption{A recursive radix-two FFT, with inverse support.}
\label{listing-fft-rec}
\end{listing}

\begin{listing}[H]
\PythonCode{listing-fft_dit}
\caption{An iterative FFT with decimation in time.}
\label{listing-fft_dit}
\end{listing}

\begin{listing}[H]
\PythonCode{listing-fft_dif}
\caption{An iterative FFT with decimation in frequency.}
\label{listing-fft_dif}
\end{listing}

\begin{listing}[H]
\PythonCode{listing-operator_s}
\caption{Multiplication by a linear phase.}
\label{listing-operator_s}
\end{listing}

\begin{listing}[H]
\PythonCode{listing-rev_bits}
\caption{Bit-reversal indices in linear total work.}
\label{listing-rev_bits}
\end{listing}

\begin{listing}[H]
\PythonCode{listing-rev_index}
\caption{Reversing the bits of a single index.}
\label{listing-rev_index}
\end{listing}

\section{Integer multiplication by FFT}
\label{sect1-listing-mult-integers}

FFT multiplication connects an approximate numerical transform with an exact
integer output. The code separates digit conversion, convolution, rounding,
and carry propagation, then checks the final product independently.

The module \texttt{polynomials.py} stores the base-$b$ digits of a
nonnegative integer from least to most significant. The conversion
functions use exact Python integers, whose precision is not fixed.
The function \texttt{multiply\_digits} pads to cover the linear
convolution, rounds its floating-point coefficients, and propagates
the carries. It rejects an excessive coefficient bound and independently
verifies the reconstructed result using Python's exact multiplication.
This verification is deliberate: rounding near an integer alone cannot
certify a floating-point computation.

The short example compares a transform product with the exact result.
For an algorithm whose transform arithmetic is itself exact, use the
modular method in Section~\ref{sect1-listing-mult-large-integers-ring}.
Neither teaching implementation is intended to improve on Python's
built-in large-integer multiplication.
\begin{listing}[H]
\PythonCode{listing-mult-integers}
\caption{FFT digit convolution and carry propagation, with exact verification.}
\label{listing-mult-integers}
\end{listing}

\begin{listing}[H]
\PythonCode{listing-number2vector}
\caption{Exact conversion from a nonnegative integer to base-$b$ digits.}
\label{listing-number2vector}
\end{listing}

\begin{listing}[H]
\PythonCode{listing-vector2number}
\caption{Reconstruction of an integer from its digits.}
\label{listing-vector2number}
\end{listing}

\begin{listing}[H]
\PythonCode{listing-test-mult-integers}
\caption{Verification of a large-integer product.}
\label{listing-test-mult-integers}
\end{listing}

\section{The Dirichlet Poisson equation}
\label{sect1-listing-fish}

The Poisson solver turns a boundary-value problem into a transform
calculation on an extended grid. Its interface makes the source, boundary
values, and grid spacing explicit so that the numerical result can be
checked against a manufactured solution.

The routines in \texttt{pde.py} implement the finite-difference problem
of Section~\ref{sect2-resol-eq-Poisson}. The full source and boundary
arrays share a grid; only the boundary entries of the second array are
used. Boundary contributions are subtracted from the interior right-hand
side, which is solved by an odd extension and a multidimensional FFT.
The helper \texttt{poisson\_zero\_boundary} includes the grid-spacing
factor in the discrete Laplacian symbol.

The manufactured solution $u(x,y)=e^{xy}$ gives
$\Delta u=(x^2+y^2)e^{xy}$ and supplies exact boundary values on the
unit square. Its left and bottom edges are constant one, while the
right and top edges are $e^y$ and $e^x$. The driver computes the grid,
numerical solution, and reference array. Refining the grid checks
second-order spatial convergence; the manufactured solution is a
verification device, not information used in the interior solve.
\begin{listing}[H]
\PythonCode{fish-listing}
\caption{A Dirichlet Poisson solver and a manufactured-solution example.}
\label{fish-listing}
\end{listing}

\begin{listing}[H]
\PythonCode{listing-f}
\caption{The source term for the exact solution $e^{xy}$.}
\label{listing-f}
\end{listing}

\begin{listing}[H]
\PythonCode{listing-sol}
\caption{An exact reference solution for the Poisson test.}
\label{listing-sol}
\end{listing}

\begin{listing}[H]
\PythonCode{listing-u_0y}
\caption{The reference boundary value at $x=0$.}
\label{listing-u_0y}
\end{listing}

\section{The periodic heat equation}
\label{sect1-listing-heat}

The periodic heat equation is especially simple in Fourier coordinates
because each mode evolves independently. The program evaluates this
evolution directly at selected times, making the effect of diffusion visible
without a time-stepping loop.

The function \texttt{heat\_periodic}, also in \texttt{pde.py}, evolves
a periodic trigonometric interpolant for
$\partial_t u=\kappa\,\partial_{xx}u$. On a period $L$, each Fourier
coefficient is multiplied by
$\exp(-\kappa t(2\pi k/L)^2)$. The grid omits the repeated endpoint.
This is spectral evolution at the requested time, rather than a
time-stepping approximation.

The driver starts from the interval indicator returned by
\texttt{heat\_initial} and returns the solutions at several times.
Changing the initial array, conductivity, period, or time does not
change the transform algorithm. Section~\ref{sect2-resol-eq-heat}
also discusses finite-difference time stepping and its separate
stability conditions.
\begin{listing}[H]
\PythonCode{heat-listing}
\caption{Heat evolution at several times from the same initial data.}
\label{heat-listing}
\end{listing}

\begin{listing}[H]
\PythonCode{listing-solve_eq}
\caption{Spectral heat evolution with an explicit period and conductivity.}
\label{listing-solve_eq}
\end{listing}

\begin{listing}[H]
\PythonCode{listing-f-heat}
\caption{An interval indicator as initial temperature.}
\label{listing-f-heat}
\end{listing}
